% !TEX root = ../main.tex
%% Corrigé du sujet type bac n° 3 — valeurs vérifiées : verifications/sujets_bac.py

\corrigeSujet{3}

\begin{corrigeBac}{8}{L'arôme de banane}
\begin{questions}
\item Le 3-méthylbutan-1-ol porte un groupe \res{\text{hydroxyle}} : c'est un \res{\text{alcool}}. L'éthanoate de 3-méthylbutyle porte un
groupe \res{\text{ester}}, $-\mathrm{COO}-$ : c'est un ester. (L'acide éthanoïque porte un groupe carboxyle.) \pts{0,75 pt}
\item $m=0{,}400\times60{,}1=\qty{24,0}{g}$, d'où $V=\dfrac{24{,}0}{1{,}05}=\res{\qty{22,9}{mL}}$. \pts{0,5 pt}
\item Le chauffage accélère la transformation ; le réfrigérant condense les vapeurs, qui retombent dans le ballon : \res{\text{pas
de perte de matière}}. L'acide sulfurique est un \res{\text{catalyseur}} : il accélère la réaction sans figurer dans son
bilan. \pts{0,75 pt}
\item L'alcool est le réactif limitant : $n_\mathrm{max}=\qty{0,200}{mol}$ d'ester, soit $0{,}200\times130{,}2=\qty{26,0}{g}$. Rendement :
$\eta=\dfrac{20{,}8}{26{,}0}=\res{80~\%}$. \pts{1,25 pt}
\item Le volume $V$ se simplifie : $K=\dfrac{x_\mathrm{f}^2}{(0{,}200-x_\mathrm{f})(0{,}400-x_\mathrm{f})}$. Avec $x_\mathrm{f}=0{,}169$ :
$\dfrac{0{,}169^2}{0{,}031\times0{,}231}=4{,}0=K$. Rendement maximal : $\dfrac{0{,}169}{0{,}200}=\res{84{,}5~\%}$. L'écart avec 80~\% vient des
pertes lors du lavage, du séchage et de la distillation. \pts{1,25 pt}
\item Avec \qty{0,200}{mol} de chaque réactif : $\dfrac{x^2}{(0{,}200-x)^2}=4$, soit $\dfrac{x}{0{,}200-x}=2$ et $x=\qty{0,133}{mol}$ : rendement de
$\res{67~\%}$. L'excès d'un réactif déplace l'équilibre vers la formation de l'ester. On choisit l'acide car il est peu coûteux
et facile à éliminer par le lavage basique, alors que l'alcool, dont la température d'ébullition (\qty{131}{\degreeCelsius}) est proche
de celle de l'ester (\qty{142}{\degreeCelsius}), serait difficile à séparer par distillation. \pts{1 pt}
\item $\ce{CH3COOH + HCO3- -> CH3COO- + CO2 + H2O}$ : le gaz est le \res{\text{dioxyde de carbone}}. Le lavage transforme l'acide
restant en ions, qui passent dans la phase aqueuse : il \res{\text{purifie}} l'ester. \pts{0,75 pt}
\item Le spectre 1 présente la large bande O--H vers \qty{3330}{cm^{-1}} : c'est celui de l'\res{\text{alcool}}. Le spectre 2 n'a pas de
bande O--H mais une bande forte vers \qty{1740}{cm^{-1}}, celle de la liaison C$=$O d'un ester : c'est celui du \res{\text{produit}}. \pts{0,75 pt}
\item L'atome d'oxygène de l'alcool, plus électronégatif que C et H, porte une charge partielle négative et des doublets non
liants : c'est le \res{\text{site donneur}}. L'atome de carbone du groupe protoné, lié à deux atomes d'oxygène plus électronégatifs
et appartenant à un ion positif, présente un déficit électronique : c'est le \res{\text{site accepteur}}. \pts{1 pt}
\end{questions}
\end{corrigeBac}

\begin{corrigeBac}{7}{Un satellite d'observation}
\begin{questions}
\item Le référentiel géocentrique a pour origine le centre de la Terre et des axes dirigés vers des étoiles lointaines.
$\vec F=\dfrac{G\,M_\mathrm{T}\,m}{r^2}\,\vec u_n$, avec $r=R_\mathrm{T}+h$ et $\vec u_n$ dirigé du satellite vers le centre de la Terre. \pts{0,75 pt}
\item $m\,\vec a=\vec F$ donne $\vec a=\dfrac{G\,M_\mathrm{T}}{r^2}\,\vec u_n$. Dans la base de Frenet, $a_t=\dfrac{\mathrm{d}v}{\mathrm{d}t}=0$ : la vitesse est constante, le
mouvement est \res{\text{uniforme}}. Et $a_n=\dfrac{v^2}{r}=\dfrac{G\,M_\mathrm{T}}{r^2}$, d'où $\res{v=\sqrt{\dfrac{G\,M_\mathrm{T}}{R_\mathrm{T}+h}}}$. \pts{1,25 pt}
\item $r=\qty{7065}{km}$ : $v=\sqrt{\dfrac{\num{6.67e-11}\times\num{5.97e24}}{\num{7.065e6}}}=\res{\qty{7,51e3}{m.s^{-1}}}$. \pts{0,5 pt}
\item $T=\dfrac{2\pi r}{v}=\dfrac{2\pi\times\num{7.065e6}}{\num{7.51e3}}=\res{\qty{5,91e3}{s}}$, soit \qty{98,5}{min} ; $\dfrac{86400}{5913}=\res{14{,}6}$ tours par jour. \pts{1 pt}
\item $T^2=\dfrac{4\pi^2r^2}{v^2}=\dfrac{4\pi^2r^3}{G\,M_\mathrm{T}}$, donc $\dfrac{T^2}{r^3}=\dfrac{4\pi^2}{G\,M_\mathrm{T}}=\qty{9,91e-14}{s^2.m^{-3}}$ : c'est la \res{\text{troisième loi de Kepler}}. \pts{0,75 pt}
\item $a=\dfrac{G\,M_\mathrm{T}}{r^2}=\res{\qty{7,98}{m.s^{-2}}}$ : la pesanteur n'est pas nulle. L'objet et le satellite tombent ensemble vers la
Terre avec la même accélération : l'objet semble flotter, c'est l'\res{\text{impesanteur}}. \pts{0,5 pt}
\item $\theta=\dfrac{\num{600e-9}}{0{,}65}=\res{\qty{9,2e-7}{rad}}$ ; au sol, à la distance $h$ : $\theta\,h=\num{9.2e-7}\times\num{694e3}=\res{\qty{0,64}{m}}$. \pts{1 pt}
\item La période vaut \qty{86164}{s} : $r=\left(\dfrac{G\,M_\mathrm{T}\,T^2}{4\pi^2}\right)^{1/3}=\qty{4,215e7}{m}$, soit une altitude de \qty{35,8e3}{km}. Le plus petit détail
visible mesurerait $\num{9.2e-7}\times\num{3.58e7}=\res{\qty{33}{m}}$ : l'orbite géostationnaire convient à la météorologie, pas à
l'observation fine du sol. \pts{1,25 pt}
\end{questions}
\end{corrigeBac}

\begin{corrigeBac}{5}{Du thé à la bonne température}
\begin{questions}
\item $\Delta U=m\,c\,\Delta T=1{,}0\times4180\times(100-18)=\res{\qty{3,43e5}{J}}$. L'eau a reçu un \res{\text{transfert thermique}} de la
résistance chauffante, elle-même chauffée par effet Joule. \pts{1 pt}
\item $\mathcal{E}=P\,\Delta t=2200\times180=\qty{3,96e5}{J}$ ; rendement $\eta=\dfrac{3{,}43}{3{,}96}=\res{87~\%}$. Pertes : échauffement de la
bouilloire elle-même, transferts thermiques vers l'air à travers les parois, évaporation d'une partie de l'eau. \pts{1 pt}
\item $\dfrac{\mathrm{d}T}{\mathrm{d}t}=-\dfrac{T_0-T_\mathrm{ext}}{\tau}\,\mathrm{e}^{-t/\tau}=-\dfrac{T-T_\mathrm{ext}}{\tau}$ : l'équation est vérifiée, et $T(0)=T_0$. \pts{1 pt}
\item À $t=\qty{10}{min}$, $T=\qty{60}{\degreeCelsius}$ : $\mathrm{e}^{-10/\tau}=\dfrac{60-20}{80-20}=\dfrac23$, d'où $\tau=\dfrac{10}{\ln1{,}5}=\res{\qty{24,7}{min}}\approx\qty{25}{min}$. \pts{0,75 pt}
\item $\dfrac{40-20}{80-20}=\dfrac13=\mathrm{e}^{-t/\tau}$, d'où $t=\tau\ln3=\res{\qty{27}{min}}$, en accord avec la courbe. Une tasse à double paroi ou
munie d'un couvercle réduit les transferts thermiques et augmente $\tau$ : le thé refroidit moins vite. \pts{1,25 pt}
\end{questions}
\schema[corrige]{bac-s3-the}
\end{corrigeBac}
